\section{Markov Mechanism Sandbox}
\label{app:v10-markov}

The Markov benchmark is the controlled sandbox for local supervision.
The oracle is known, the sufficient state is known, and node labels
can be assigned at designed propensities. This isolates the mechanism
that the Benoit replication leaves mixed with measurement noise:
local labels can train and audit the state map when they measure the
same compositional property as the root label.

\subsection{DGP and Sufficient State}

Each synthetic document is a fixed-length token sequence generated by
a hidden register process. The learner sees observed tokens from
register-specific synonym palettes. The oracle is the number of
hidden-register transitions in the sequence.

For any span, the sufficient state is
\[
  (\mathrm{count},\mathrm{first},\mathrm{last}).
\]
The count records transitions inside the span. The first and last
fields carry the boundary information used by parent merges. The
merge rule is
\[
(c_L,a_L,b_L)\otimes(c_R,a_R,b_R)
=
(c_L+c_R+\One\{b_L\ne a_R\},a_L,b_R).
\]
The state is ordered and associative. A parent merge needs the right
boundary of the left child and the left boundary of the right child;
dropping those fields creates exactly the kind of C3 failure that
C-TreePO audits. The point of the benchmark is therefore not that
Markov counts are hard. It is that local supervision has to recover
the state needed for future composition, rather than only the local
scalar that looks predictive in isolation.

\begin{figure}[t]
\centering
\includegraphics[width=\linewidth]{markov_scene_arc_hamlet.pdf}
\caption{\textbf{Boundary correction in the Markov sandbox.}
Twelve consecutive scenes are used as a reader-facing skin on the
synthetic registers. Each leaf stores
\((\mathrm{count},\mathrm{first},\mathrm{last})\). Internal nodes add
the child counts and a boundary correction when adjacent child
registers differ.}
\label{fig:v10-markov-registers}
\end{figure}

\subsection{Leaf-Size and Allocation Grid}

The completed grid uses \(10{,}240\) training documents, leaf sizes
\(128,64,32,16,8\), and root-label shares from \(100\%\) down to
\(10\%\). The recoverable regime has \(4\) hidden registers and about
\(5\) expected changepoints per document. The structural regime has
\(12\) hidden registers and about \(10\) expected changepoints per
document. Both regimes use disjoint observed-token palettes so the
root task is recoverable from the document, and remaining error is
about representation, local supervision, and optimization.

\begin{figure}[t]
\centering
\begin{subfigure}[b]{0.48\textwidth}
    \centering
    \includegraphics[width=\linewidth]{markov_leaf_size_fixed_recoverable.pdf}
    \caption{Recoverable regime.}
    \label{fig:v10-markov-leaf-size-recoverable}
\end{subfigure}
\hfill
\begin{subfigure}[b]{0.48\textwidth}
    \centering
    \includegraphics[width=\linewidth]{markov_leaf_size_fixed_structural.pdf}
    \caption{Structural regime.}
    \label{fig:v10-markov-leaf-size-structural}
\end{subfigure}
\caption{\textbf{Tree-supervision gain over the flat FNO baseline
across leaf size and supervision budget.} The recoverable regime has
a broad interior band where tree supervision matches or improves on
the flat baseline. The structural regime is harder and exposes the
boundary-state failure mode at small leaves.}
\label{fig:v10-markov-leaf-size}
\end{figure}

Figure~\ref{fig:v10-markov-leaf-size} reads the state-size
tradeoff. Very large leaves make the tree close to a flat root
reader. Very small leaves increase the number of boundary decisions.
The useful region is the interior band where the state carries enough
boundary information and the local labels give enough training signal.

\begin{figure}[t]
    \centering
    \includegraphics[width=0.92\linewidth]{markov_budget_split.png}
    \caption{\textbf{Splitting a fixed token budget between root and
    local labels.} The green star is the all-root reference. The
    green line spends only the reduced root-label budget. The blue,
    purple, and rose curves keep total supervision mass fixed by
    assigning the missing mass to local labels under three allocation
    policies.}
    \label{fig:v10-markov-budget-split}
\end{figure}

Figure~\ref{fig:v10-markov-budget-split} is the supervision-granularity
intuition check. A full-document label on a \(128\)-token sequence costs
\(128\) token observations. A label on a \(16\)-token leaf costs
\(16\) token observations, so eight such leaf labels have the same
full-document-equivalent mass as one root label. Reallocating missing
root-label mass to local labels keeps root MAE close to the all-root
reference over a substantial range of root-label shares.

\subsection{Provenance}

The allocation-policy grid is complete for the planned \(100\%\) to
\(10\%\) root-label shares: \(123\) scheduled items completed with
\(0\) failures in
\path|outputs/markov_v5_simple_fixed10240_quick_20260414_utc/combined_scheduler_allocation_policy_grid/|.
The rendered paper figures are in
\path|paper/ctreepo/assets/markov/figures/| and
\path|outputs/markov_v5_sticky_allocation_policy_paper_20260416/|.
The generated coverage report records full replacement coverage for
both regimes at root shares \(90\%\) through \(10\%\), including the
leaf-only, depth-equal, and balanced-node allocation policies where
the tree has non-root nodes to distinguish them.
